\documentclass[a4paper,11pt,twoside]{article}
\usepackage[utf8]{inputenc}	% Text coding
\usepackage[T1]{fontenc}
\usepackage{lmodern}
\usepackage[czech]{babel}
\usepackage{epsfig}
\usepackage{amsfonts,amsmath,amssymb}
\usepackage{graphicx}
\usepackage[unicode]{hyperref}
%\usepackage{indentfirst}
\usepackage{fancyhdr}
\usepackage{xifthen}
\usepackage{amsthm,thmtools}
\usepackage{bold-extra}
\usepackage[dvipsnames]{xcolor}
\usepackage[subrefformat=simple,labelformat=simple]{subcaption} % Instead of subfigure
\usepackage{listings}
\usepackage{comment}
\usepackage{titlesec}
\usepackage{underscore}
\usepackage{makecell}       % Šířky čar v tabulkách

% Page size
\addtolength{\topmargin}{-1.5cm} %\addtolength{\textheight}{-10cm}
\addtolength{\textwidth}{4cm} \addtolength{\textheight}{4cm} % Width and height of the text
\addtolength{\voffset}{-0.5cm} % Top margin
\addtolength{\hoffset}{-2cm}
\setlength{\headheight}{15pt}

\DeclareMathOperator{\e}{e}

\def\vector#1{\boldsymbol{#1}}								% Vector
\renewcommand{\d}{\mathrm{d}}
\newcommand{\derivative}[3][]{\ifthenelse{\isempty{#1}}	    % Normal derivative
	{\frac{\d{#2}}{\d{#3}}}
	{\frac{\d^{#1}{#2}}{\d{#3}^{#1}}}
}
\newcommand{\im}{\mathrm{i}}

\def\makematrix#1{\begin{pmatrix}#1\end{pmatrix}}       % Matrix
\def\abs#1{\left|#1\right|}
\def\probability#1{\mathrm{Pr}\left[#1\right]}
\def\expectation#1{\mathrm{E}\left[#1\right]}
\def\dispersion#1{\sigma_{#1}^{2}}
\def\c{,\!}

\def\code#1{\textnormal{\texttt{#1}}}
\def\file#1{\textnormal{\textbf{\texttt{#1}}}}
\def\ghfile#1#2{\textnormal{\textbf{\texttt{\href{https://github.com/PavelStransky/PCInPhysics/blob/main/#1#2}{#2}}}}}

\def\abbreviation#1{\textnormal{\textsc{#1}}}

\def\ket#1{\left|{#1}\right\rangle}
\def\bra#1{\left\langle{#1}\right|}
\def\braket#1#2{\left\langle{#1}\middle|{#2}\right\rangle}

\begin{document}

\section*{Domácí úkol na 19.5.2025}
\subsection*{Diskrétní Fourierova transformace}

\begin{enumerate}
    \item Naprogramujte Fourierovu transformaci, zpětnou Fourierovu transformaci pro zadanou časovou řadu $h_j$, $j=0,1,\dotsc,N-1$ a funkci pro výpočet výkonového spektra $S_k$.

    \item
        Vytvořte časovou řadu délky $N=2000$ se vzorkovací frekvencí $f_{s}=2000\,\text{Hz}$ (signál tedy bude trvat přesně $1\,\text{s}$) danou součtem tří harmonických funkcí:
        \begin{equation*}
            h_{j}=\sum_{n=1}^{3}a_{n}\sin\left(2\pi F_{n} t_{j}\right),
            \label{eq:FTTest}
        \end{equation*}
        kde
        \begin{align}
            a_{1}&=0.1, & a_{2}&=0.2 & a_{3}&=0.3,\nonumber\\
            F_{1}&=440\,\text{Hz} & F_{2}&=\frac{5}{4} F_{1}=550\,\text{Hz} &F_{3}&=\frac{3}{2}F_{1}=660\,\text{Hz}\nonumber
        \end{align}
        a časy $t_{j}=j/f_{s}$, $j=0,1,\dotsc,N-1$.
        Spočítejte Fourierovu transformaci a vykreslete graf výkonového spektra.
        Ověřte, že graf má tři vrcholy s výškami odpovídajícími zadaným amplitudám $a_{n}$ a středy v místech zadaných frekvencí $F_{n}$.

    \item 
        Spočítejte Fourierovu transformaci pro časovou řadu $\theta_{j}$, kterou získáte navzorkováním řešení $\theta(t)$ pohybové rovnice pro buzené a tlumené kyvadlo z \href{https://github.com/PavelStransky/PCInPhysics/blob/main/ukoly/2025/Ukol5.pdf}{Úkolu 5} s parametry $\eta=0.5\,\text{Hz}, \omega^2=A=1\,\mathrm{s}^{-2}$ na časovém intervalu $t\in\langle 0\,\mathrm{s},1000\,\mathrm{s}\rangle$ a s vzorkovací frekvencí $f_{s}=2\,\text{Hz}$.
    
        Volte tři různé frekvence buzení $\omega_{b}=0.2\,\mathrm{s}^{-1}, 0.5\,\mathrm{s}^{-1}, 0.8\,\mathrm{s}^{-1}$ a pro každou z nich spočítejte Fourierovu transformaci a vykreslete graf výkonového spektra.
        V případě $\omega_{b}=0.2\,\mathrm{s}^{-1}$ a $\omega_{b}=0.8\,\mathrm{s}^{-1}$ porovnejte dominantní frekvenci Fourierovy transformace s frekvencí buzení $\omega_{b}$.
        V chaotickém případě $\omega_{b}=0.5\,\mathrm{s}^{-1}$ dostanete \uv{spojité} spektrum frekvencí.
        Vykreslete ji v log-log grafu. 
        Proložte přímku 
        \begin{equation*}
            y=ax+b
        \end{equation*}
        hodnotami $x_{k}=\log{f}_{k},y_{k}=\log{S}_{k}$, kde $S_{k}$ je hodnota výkonového spektra pro frekvence $f_{k}$ omezené $f_{k}<200\,\mathrm{Hz}$ a určete sklon přímky $a$.
        K prokládání přímky použijte například\\
        \code{a,b = numpy.polyfit(x, y, deg=1)}.

\end{enumerate}


\end{document}